\documentclass[12pt,a4paper]{article}

\usepackage[utf8]{inputenc}
\usepackage[T1]{fontenc}
\usepackage{times}
\usepackage[margin=1in]{geometry}
\usepackage{setspace}
\usepackage[round,authoryear]{natbib}
\usepackage[colorlinks=true,linkcolor=blue,citecolor=blue,urlcolor=blue]{hyperref}
\usepackage{microtype}
\usepackage{abstract}

\setlength{\parindent}{0pt}
\setlength{\parskip}{6pt}

\onehalfspacing

\begin{document}


\begin{center}
{\LARGE\bfseries The B+ Trap}\\[0.5em]
{\Large How AI Compresses the Creative Spectrum ---}\\[0.3em]
{\Large And What a Rebel Architecture Could Change}\\[2em]
{\large Fabio Lauria}\\[0.3em]
{CEO \& Founder, ELECTE}\\
{Milan, Italy}\\[1.5em]
{February 2026}\\[0.5em]
{\small PREPRINT}
\end{center}

\vspace{1em}

\begin{abstract}
\noindent
Generative artificial intelligence now surpasses average human performance on standard creativity benchmarks, yet the most creative humans still outperform every AI system tested. This asymmetry defines what we call the B+ Trap: AI compresses the creative spectrum from both directions, pulling weaker performers up while thinning the top of the distribution as fewer creators push beyond AI's suggestions. A 2025 NeurIPS Best Paper studying 70+ state-of-the-art language models confirmed the effect at ecosystem scale, documenting an ``Artificial Hivemind'' of extreme inter-model homogenization. A meta-analysis of 28 studies (8,214 participants; \citealt{holzner2025}) quantifies the trade-off---human-AI collaboration yields a moderate boost to individual creative performance (Hedges' g = 0.27) but inflicts a large negative effect on idea diversity (g = $-$0.86). The mechanism is structural: Reinforcement Learning from Human Feedback (RLHF) provably narrows output distributions and amplifies majority preferences, and alignment training has been shown to systematically reduce language models' conceptual diversity. Yet some creators consistently escape the trap. The research identifies a key differentiator: metacognitive engagement---the capacity to critically evaluate, actively diverge from, and strategically redirect AI suggestions rather than accept them. This paper argues that what these exceptional humans do cognitively, a different class of AI architecture could do computationally. We present a speculative but empirically grounded ``Rebel AI'' framework---synthesizing adversarial creativity networks, novelty-maximizing reward functions, curiosity-driven exploration, and open-ended evolutionary search---as alternative training paradigms that optimize for divergence rather than convergence. We conclude with implications for small and medium enterprises, where creative distinctiveness is often the primary competitive advantage.


\noindent \textbf{Keywords: artificial intelligence, creativity, RLHF, homogenization, adversarial training, novelty search, metacognition, SME strategy, divergent thinking, generative AI}
\end{abstract}

\section{Introduction}
A curious paradox is emerging across the AI-assisted creative landscape. By every standard measure, generative AI makes individuals more creative. Stories written with AI assistance are rated more enjoyable. Marketing copy produced faster. Design iterations generated in seconds rather than days. Yet something is being lost in the aggregate---and something specific is being lost at the top. The outputs are converging. The creative marketplace is flattening. The floor is rising, but the ceiling is dropping. Everyone is becoming a B+ creative.

We call this the B+ Trap, and it operates through a double compression. From below, AI lifts weaker performers toward competence---less creative writers produce stories rated up to 26.6\% better when assisted by GPT-4 \citep{doshi2024}. From above, AI does not make top performers individually worse, but it thins the top of the distribution: by anchoring everyone's creative direction to the same probability distributions, it reduces the number of people who attempt to push beyond the AI's suggestions. The trap is that B+ feels good enough. The output is polished, fluent, and defensible. The motivation to push further dissolves.

Yet not everyone gets trapped. The largest comparative study on AI and human creativity \citep{bellemare2026} found that the top 10\% of human scorers still exceed every AI model tested, and the gap widens at the extremes. \citet{wang2025} in Nature Human Behaviour, comparing 9,198 humans against 215,542 LLM observations, confirmed that human creativity is most pronounced at the extremes---humans exhibit greater variability and higher peak creativity than any AI system. The question that animates this paper is: what differentiates these people? Across multiple independent studies, a key differentiator is metacognitive engagement---the capacity to critically evaluate AI outputs, actively diverge from them, and use them as departure points rather than destinations. The people who escape the B+ Trap are, in essence, rebelling against the AI's gravitational pull toward the probable.

This observation leads to the paper's central argument, which operates at two levels. At the human level, we synthesize the rapidly expanding 2024--2026 research documenting the B+ Trap's mechanics---how RLHF structurally causes creative narrowing, how the effect compounds at population scale, and why metacognition is the escape mechanism. At the architectural level, we argue that what exceptional humans do cognitively, a different class of AI could do computationally. We present a speculative framework---what we term ``Rebel AI''---that synthesizes adversarial creativity networks, novelty-maximizing reward functions, curiosity-driven exploration, and open-ended evolutionary search into a counter-paradigm: training AI to deviate from, rather than converge on, learned patterns.

The stakes are particularly acute for small and medium enterprises (SMEs). Unlike large corporations with established brand equity, SMEs often compete on distinctiveness, personality, and the capacity for creative differentiation. A natural experiment during Italy's 2023 ChatGPT ban provides striking evidence: when Milan restaurants lost AI access, their Instagram content showed 15\% less lexical similarity---and generated approximately 3.5\% more engagement \citep{liu2025italy}. The homogenized AI content was not merely similar; it was less effective at driving consumer response.

The paper proceeds as follows. Section 2 reviews the empirical evidence on AI creative performance relative to humans. Section 3 establishes the homogenization paradox. Section 4 identifies RLHF as the structural mechanism. Section 5 examines metacognition as a key differentiator separating those who escape the B+ Trap from those who remain in it. Section 6 situates these findings within computational creativity theory. Section 7 presents the Rebel AI framework. Section 8 discusses AlphaEvolve as a proof-of-concept. Section 9 draws implications for SME deployment strategy. Section 10 acknowledges limitations. Section 11 concludes.

\section{The Empirical Record: AI Surpasses the Average, Not the Best}
The empirical record from 2024--2026 is now extensive enough to draw confident conclusions about where AI creativity stands relative to human creativity. What emerges is not a simple story of AI superiority or inferiority, but a consistent pattern: AI outperforms the median human while falling short of the top performers, creative performance across model generations may have plateaued, and the gap at the extremes shows no sign of closing with model scale.

\subsection{Large-Scale Comparative Studies}
The largest comparative study ever conducted---\citet{bellemare2026} in Scientific Reports, involving over 100,000 human participants and seven major LLMs including GPT-4, Claude 3, and Gemini Pro---found that GPT-4 surpassed average human scores on the Divergent Association Task (DAT) with statistical significance. GeminiPro performed indistinguishably from the human average. Yet the top 50\% of human scorers exceeded every AI model tested, and the gap widened further at the top 10th percentile. This study, authored by a team including Yoshua Bengio at Mila, also revealed a striking word repetition problem: GPT-4-turbo used ``ocean'' in over 90\% of its responses, while the most common human word (``car'') appeared in only 1.4\% of responses.

\citet{wang2025} in Nature Human Behaviour provided the most statistically powerful confirmation, comparing 9,198 human participants against 215,542 LLM observations from nine models on the DAT. Human creativity on average was slightly higher than LLMs, with differences most pronounced at the extremes---humans exhibit greater variability and higher peak creativity. Increasing the ``creativity'' temperature parameter improved LLM performance up to a threshold, then outputs became nonsensical. Prompt engineering yielded mixed to negative results: genius persona prompting and demographic role-playing reversed real-world creativity patterns rather than enhancing performance.

\citet{hubert2024} in Scientific Reports tested 151 humans against GPT-4 across a full battery of divergent thinking tasks---the Alternate Uses Task (AUT), Consequences Task, and DAT---finding GPT-4 was robustly more creative across all three in originality and elaboration, even after controlling for fluency. \citet{guzik2023} in the Journal of Creativity administered the Torrance Tests of Creative Thinking to GPT-4, which scored in the top 1st percentile nationally for fluency and originality against 2,700 college student norms. \citet{ruan2025} conducted the largest-scale creativity comparison of models, evaluating 41+ systems including o1-preview, o3-mini, Claude 3.7 Sonnet, GPT-4.5 Preview, and DeepSeek-R1 across 1,180 keywords spanning 22 scientific domains. Scientific creative ability showed distinct patterns from general intelligence metrics---high benchmark scores did not predict creativity---and a clear Pareto frontier emerged between feasibility and originality across models.

\subsection{The Ceiling Effect and the Creativity Plateau}
The ceiling effect holds consistently across studies. \citet{koivisto2023}, testing 256 humans against multiple AI systems on the AUT, found that while AI chatbots outperformed humans on average semantic distance and creativity ratings, the best human ideas matched or exceeded chatbot outputs in seven of eight scoring comparisons. The Bellemare-Pepin study confirmed this pattern held even with frontier models: creative writing tasks showed LLMs approaching but not reaching the performance of highly creative humans. \citet{rosenbaum2025} tested GPT-3.5, GPT-4, Claude 1.3, Qwen, and SparkDesk across 13 creative tasks and found the best LLMs ranked at only the 52nd percentile against individual humans, excelling in divergent thinking but lagging consistently in creative writing.

Perhaps the most striking finding comes from \citet{haase2025}, published in the Journal of Creativity. Testing 14 LLMs including GPT-4o and o3-mini on both the DAT and AUT with standard and ``creative'' prompts, they found no evidence of increased creative performance over 18--24 months of model development. GPT-4 actually performed worse than in prior studies. On the AUT, only 0.28\% of LLM-generated responses reached the top 10\% of human creativity benchmarks---humans are roughly 35.7 times more likely to produce standout creative ideas. This plateau persists despite dramatic improvements on reasoning benchmarks during the same period, suggesting that creative capacity and general intelligence follow different scaling trajectories.

Critically, newer model versions do not always improve creative performance---GPT-4-turbo showed a notable decline relative to GPT-4, demonstrating that capability scaling does not automatically translate to creative gains. The CREATIVITY INDEX developed by Lu et al. (ICLR 2025, oral, top 1.8\% of submissions) provides the most rigorous quantification: professional human authors' creativity index is 66.2\% higher than LLMs on average, positioning LLMs as ``remixers'' rather than original creators. The index measures how much text can be reconstructed from existing web text snippets---a direct operationalization of originality. Lu et al.'s RLHF-specific finding is equally important: alignment training reduces LLM creativity by approximately 30.1\%, suggesting that the very process intended to make models more useful also makes them less original.

\citet{zhang2025creative} reinforced the distinction between output quality and creative process through eight pre-registered experiments comparing GPT-3.5, GPT-4, and GPT-4o against human participants across AUT, DAT, free association, insight problem-solving, creative writing, and creative evaluation. AI outperformed humans on structured divergent thinking tasks but underperformed in creative writing, showed higher response repetition rates, and exhibited no significant correlation between novelty and appropriateness weights---suggesting the absence of the human-like trade-off strategy that underlies genuine creativity.

\subsection{The Diversity Deficit and the Model Generation Question}
NoveltyBench \citep{zhang2025creative} exposes a systematic blind spot in how the field evaluates AI creativity. Using 1,100 prompts measuring LLM ability to generate multiple distinct, high-quality outputs via functional equivalence classes, the authors found that all 20 frontier models tested---including GPT-4o, Claude 3.5 Sonnet, and Gemini 2.0 Pro---suffer from lack of diversity. A sobering structural finding: larger, more capable models within a family produce less diverse outputs. Over 90\% of benchmark papers at COLM 2024 and ICLR 2025 evaluate models on single or best generation only, systematically ignoring the diversity dimension. The field, in other words, is optimizing for a metric (peak quality) that obscures the problem (distributional narrowing). \citet{shypula2025} at ICLR's DL4C Workshop formalized this through the concept of ``effective semantic diversity''---diversity among outputs meeting quality thresholds---and found that preference-tuned models exhibited reduced lexical diversity but preserved semantic diversity, a distinction traditional metrics overlook.

A note on model generations is warranted. The creativity research reviewed in Sections 2--5 was conducted predominantly with GPT-4 era models (GPT-3.5, GPT-4, GPT-4o, Claude 3.5 Sonnet, Gemini 2.0). As of early 2026, frontier models have advanced considerably---GPT-5.2, Claude Opus 4.6, Gemini 3 Pro. Three sources now illuminate how creativity behaves across model generations.

First, the Creativity Benchmark \citep{bhat2025}, developed with seven advertising organizations and evaluated by 678 practicing creatives through 11,012 pairwise comparisons, tested 14 models including O3, Gemini 2.5 Pro Preview, DeepSeek Reasoner, Claude 3.7 Sonnet, LLaMA 4 Scout, and GPT-4.5 Preview on TTCT/Guilford divergent thinking dimensions. O3 attained the highest originality scores. Yet overall performance was tightly clustered ($\Delta\theta$ $\approx$ 0.45), and correlations between LLM-judge originality scores and human preferences were weak (Spearman $\rho$ $\approx$ 0.32)---models are converging on similar creative output even as they advance on other dimensions.

Second, Haase et al.'s plateau finding is particularly significant because it includes o3-mini---a reasoning-optimized frontier model---and still finds no creative improvement over 18--24 months. If reasoning-chain enhancements do not translate to creative gains, the B+ Trap may persist even as models advance on other cognitive dimensions.

Third, \citet{wright2025} tested 27 LLMs including GPT-5 across 155 topics and 200 prompt variations, measuring epistemic diversity---the variation in real-world claims across outputs. Newer models, particularly GPT-5, show sharp increases in epistemic diversity compared to predecessors. However, nearly all models tested remain less epistemically diverse than a basic web search, and model size has a negative impact on diversity within model families. These results are consistent with the theoretical work reviewed in Section 4---\citet{gxchen2025} and J. \citet{chen2025entropy} demonstrate that mode collapse is a mathematical property of standard KL-regularized RL objectives regardless of model scale.

Taken together, the frontier model evidence suggests the B+ Trap is softening but not dissolving: newer models show marginal improvements on some diversity dimensions, but the structural gap between AI and human-level creative diversity persists, creative performance may have plateaued, and the most advanced reasoning models do not escape the pattern. No published creativity benchmark study has yet tested GPT-5.2, Claude Opus 4.6, or Gemini 3 Pro on standardized instruments such as the AUT, TTCT, or creative writing evaluation---this remains an open empirical question. Until frontier models adopt fundamentally different training paradigms---of the kind described in Section 7---the B+ Trap's structural logic should hold, even as its severity may gradually diminish.

\section{The Homogenization Paradox}
The NeurIPS 2025 Datasets and Benchmarks Best Paper---\citet{jiang2025}---provided the most authoritative demonstration of what they termed the ``Artificial Hivemind.'' Evaluating over 70 state-of-the-art LLMs using INFINITY-CHAT, a dataset of 26,000 diverse real-world open-ended queries with 31,250 dense human annotations (25 independent ratings per prompt), they documented extreme intra-model repetition and inter-model homogenization: different models from different companies produce strikingly similar outputs. Existing reward models and automated evaluators were poorly calibrated to the diversity of actual human judgments. This is the B+ Trap demonstrated at ecosystem scale---not a property of any single model, but an emergent characteristic of the entire LLM landscape.

The landmark study by \citet{doshi2024}, published in Science Advances, crystallized what has become the defining empirical finding. In a preregistered experiment with approximately 300 participants writing short stories under three conditions---no AI, one GPT-4 idea, or five GPT-4 ideas---AI-assisted stories were rated as more creative, better written, and more enjoyable. Less creative writers benefited most, with their stories rated up to 26.6\% better written and 15.2\% less boring. But AI-enabled stories were 5.2\% more similar to each other (b = 4.29, P < 0.001) than human-only stories. The authors described this as a ``social dilemma'': individually rational adoption of AI produces collectively suboptimal outcomes for creative diversity.

\citet{wenger2025}, working across Duke University and the Technion, tested 22 LLMs using standardized creativity tests---AUT, Forward Flow, and DAT---comparing population-level diversity against human baselines. The quantitative gap is stark: mean cosine distance variability was 0.459 for LLMs versus 0.738 for humans on the AUT. T-SNE visualization showed dense clustering of LLM outputs versus dispersed human outputs. The critical finding: homogeneity persisted even after controlling for response structure, verbosity, model size, and vendor. Creative homogeneity is not limited to individual models---it is systemic across all LLMs tested, suggesting shared architectural and training biases produce converging creative outputs regardless of vendor differentiation.

\citet{anderson2024}, presented at the ACM Creativity \& Cognition conference, provided the mechanistic detail. In a 36-participant within-subjects study generating 1,271 ideas across product design and fictional scenario tasks, ideas produced with ChatGPT showed significantly less divergence from the group average (M = 0.24 cosine similarity) compared to those produced with the non-AI tool Oblique Strategies. The critical finding was that homogenization operated at the group level, not the individual level---ChatGPT did not make any single person's ideas more repetitive, but it pushed different people's ideas toward the same center of gravity. Participants reported that ChatGPT made them ``turn my brain off,'' while the analog tool was ``harder to use, but it was more fun and it got me to more interesting places.''

This effect scales badly. \citet{moon2025} analyzed 2,200 college admissions essays using a novel ``diversity growth rate'' metric and found that each additional human-written essay contributed 2--8 times more new ideas to the collective pool than a GPT-4 essay, with the homogenization gap widening as more essays were added. Even parameter and prompt modifications that improved individual-level diversity did not close the collective diversity gap, demonstrating that the problem is architectural rather than configurable. Xie, Jurafsky, Riedl, et al. (2025) in PNAS introduced the Sui Generis score---an automatic metric measuring uniqueness of plot element combinations---and found that LLM-generated stories contained repetitive combinations of idiosyncratic plot elements echoed frequently across generations and across different models, while human-written stories maintained significantly higher plot uniqueness.

The depth of homogenization extends beyond lexical patterns. \citet{alvero2025} analyzed 372,793 U.S. college admissions essays and found a ``paradoxical homogenization'': post-ChatGPT essays showed greater lexical diversity but lower idea diversity---surface-level word variety masks deeper conceptual convergence. This pattern inflated creativity ratings because diverse vocabulary was perceived as more creative despite fewer unique ideas. Effects were particularly pronounced among racial and linguistic minority applicants. An information-theoretic analysis \citep{sui2026} comparing LLM creative fiction against professional writers (MFA graduates, published authors) found that LLMs show systematically lower entropy, with RLHF exacerbating mode collapse and KL-divergence regularization in alignment penalizing the long tail of diverse outputs.

The cultural dimension is equally concerning. \citet{agarwal2025}, presented at CHI, ran a cross-cultural controlled experiment with 118 participants from India and the US on culturally grounded writing tasks. AI suggestions led Indian participants to adopt Western writing styles, altering both content and form---the first empirical evidence of cultural homogenization from embedded AI writing features. \citet{sourati2025} synthesized evidence across linguistics, psychology, cognitive science, and computer science, concluding that next-token prediction fundamentally favors high-probability continuations, smoothing over outliers, with LLM-authored content disproportionately reflecting affluent subgroup writing styles and warning of ``flattening the cognitive landscapes that drive collective intelligence.''

Most alarmingly, \citet{zhou2026scar} in Technology in Society found in a longitudinal study that creativity drops significantly upon AI withdrawal, and induced content homogeneity continues climbing even months after AI use stops, creating what they term a ``creative scar.'' The implications are profound: the homogenization effect is not merely a property of AI-generated content but becomes embedded in the cognitive habits of users who rely on it. A meta-analysis by \citet{holzner2025} of 28 studies (8,214 participants) quantifies the overall trade-off: human-AI collaboration yields a moderate boost to individual creative performance (Hedges' g = 0.27) but inflicts a large negative effect on idea diversity (g = $-$0.86).

\citet{liu2025italy} leveraged Italy's country-wide ChatGPT ban in April 2023 as a natural experiment, studying restaurant Instagram content. During the ban, restaurants showed relative decreases of 15\% in lexical similarity and 12\% in syntactic similarity, alongside a 3.5\% increase in average consumer engagement---the most direct causal evidence yet that AI access drives content homogenization and that homogenized content is less commercially effective.

\section{RLHF: The Structural Mechanism of Creative Narrowing}
The empirical literature now identifies Reinforcement Learning from Human Feedback (RLHF) as a primary driver of the diversity collapse observed in creative AI outputs. \citet{padmakumar2024}, presented at ICLR 2024, ran a controlled experiment comparing writing with base GPT-3, InstructGPT (RLHF-tuned), and no model assistance. Writing with InstructGPT produced a statistically significant reduction in content diversity, but base GPT-3 did not---directly implicating the alignment process rather than pretraining as the cause. The diversity reduction came entirely from the model's contributions; human-written portions remained unaffected.

\citet{kirk2024} at ICLR provided the first rigorous empirical demonstration that RLHF causes both per-input and across-input mode collapse. Comparing SFT to RLHF across LLaMA-7B and OPT models, they found RLHF substantially decreased output diversity on every metric tested (self-BLEU, semantic similarity, NLI). \citet{mohammadi2024} showed that RLHF-aligned Llama-2 models exhibit lower token-level entropy, form distinct embedding-space clusters, and gravitate toward ``attractor states''---recurring output patterns the model converges to regardless of prompt variation. \citet{hamilton2024} at the SCALE-LLM workshop documented that successive GPT-3 versions, as alignment training increased, showed increasing degrees of mode collapse in narrative generation, losing the ability to generate stylistically distinct authorial voices.

\citet{krishnamurthy2025}, published at NAACL, provided the mechanistic explanation. Their study demonstrated that alignment---including both RLHF and instruction tuning---systematically reduces language models' conceptual diversity. This is not a side effect but a direct consequence of how alignment operates: by training models to converge on human-preferred outputs, the process eliminates the conceptual outliers that constitute the long tail of creative possibility. \citet{guo2025} in Transactions of the ACL developed a comprehensive framework evaluating lexical, syntactic, and semantic diversity across LLMs, finding that instruction tuning improves lexical diversity but constrains syntactic and semantic diversity---a narrowing of expressive flexibility beneath surface-level word variety that mirrors the paradoxical homogenization observed by Alvero et al. in human writing under AI influence.

The theoretical foundation is now established. \citet{xiao2025}, accepted in the Journal of the American Statistical Association, proved mathematically that RLHF's KL-divergence regularization contains an inherent algorithmic bias that asymptotically amplifies preference imbalances---minority preferences are virtually eliminated in the limit. This ``preference collapse'' persists even when the reward model is a perfect oracle. Their proposed fix, Preference Matching RLHF, achieved 29--41\% improvement in alignment with actual human preference distributions. \citet{xu2025} in PNAS quantified the narrative-level consequences: when given a Kafka story prompt, 50 of 100 GPT-4 continuations had the policeman give instructions to ``take the second left,'' and 16 mentioned a bakery as a landmark. None resembled Kafka's actual ending.

The probability-novelty tradeoff is mathematically fundamental. \citet{holtzman2020} established that maximization-based decoding produces degenerate text because even the best language models have unreliable probability tails. Welleck et al. (ICLR 2020) showed the maximum likelihood objective itself assigns excessive probability to repeated sequences. \citet{gxchen2025} provide theoretical analysis in their MARA framework showing that standard reverse-KL regularized RL objectives inherently promote mode collapse because small reward differences drive exponentially large probability shifts. J. \citet{chen2025entropy} formalizes model collapse across LLMs, GANs, and RL as entropy decay in self-referential learning systems via the Entropy-Reservoir Bregman Projection framework, proving that coupling to external entropy sources is both necessary and sufficient for stability. Hintze et al. (\emph{Patterns}, 2025) provide the most vivid demonstration: autonomous AI image-text loops consistently converge to just 12 generic motifs regardless of initial prompt diversity. The creative gradient, without intervention, always flows downhill.

\section{Who Escapes the B+ Trap: Metacognition as the Differentiator}
If RLHF is the structural mechanism of the B+ Trap, metacognition is a key factor separating those who escape from those who remain in it. The trap operates through cognitive disengagement---the seductive comfort of polished, good-enough output that dissolves the motivation to push further. The people who reach A+ are not using different tools; they are using the same tools with a fundamentally different cognitive posture.

\citet{kumar2025} at CHI conducted the most comprehensive experimental demonstration, running two pre-registered parallel experiments with 1,100 participants in three conditions: standard LLM, coach-like LLM, and no assistance. Using the AUT for divergent thinking and RAT for convergent thinking, with cosine distance measured via SBERT embeddings across 150 Monte Carlo runs, they found that LLM-generated strategies (coach mode) reduced diversity both during and after LLM use---a lasting impact on creative processes. Participants with no LLM interactions generated the most diverse idea sets collectively. LLM guidance led to decreased originality (from 2.77 to 2.69) in subsequent unassisted tasks, demonstrating that the B+ Trap's effects persist beyond the moment of AI interaction.

\citet{fernandes2026} in Computers in Human Behavior conducted the most direct test of metacognitive calibration: across two studies (N = 246 and N = 452), participants using ChatGPT for logical reasoning improved their performance by 3 points relative to population norms but overestimated their performance by 4 points. Higher AI literacy correlated with lower metacognitive accuracy---technically sophisticated users were more confident but less precise in self-assessment. The Dunning-Kruger effect disappeared entirely under AI assistance, replaced by a universal metacognitive blind spot.

\citet{lee2025} at CHI surveyed knowledge workers and found that confidence in AI's ability to perform tasks showed a large negative association with critical thinking engagement ($\beta$ = $-$0.69, p < 0.001), while confidence in one's own ability ($\beta$ = 0.26) and confidence in evaluating AI outputs ($\beta$ = 0.31) both positively predicted critical thinking. This is the mechanism of escape: those who trust their own judgment over the AI's, and who feel equipped to evaluate rather than simply accept, are the ones who diverge productively.

\citet{maier2025} provided the most nuanced experimental evidence in randomized controlled experiments with 486 and 640 participants using the AUT across five conditions. Model-led interaction substantially improved idea quality but reduced idea diversity and perceived ownership. The reflective, human-led question mode improved quality while preserving diversity and ownership---a direct experimental demonstration that how humans engage with AI matters more than whether they engage. When the human leads, the B+ Trap opens; when the AI leads, it closes.

The timing of AI engagement matters as well. A controlled experiment at CHI 2025 with 60 participants on health-related ideation \citep{qin2025} found that using LLMs from the beginning led to idea fixation, reduced creative self-efficacy, and lowered self-credit. Delaying LLM use until after independent ideation increased quantity and decreased similarity to LLM suggestions, with autonomy and ownership fully mediating the improvement. This aligns with the Doshi and Hauser data: less creative writers benefited most from AI but showed the strongest anchoring to AI-generated ideas, converging on suggestions rather than building beyond them.

\citet{wadinambiarachchi2024} at CHI demonstrated the mechanism through a design experiment: 60 designers using Midjourney produced fewer ideas, with less variety and lower originality compared to those working without AI, due to design fixation. The critical variable was not AI access but how users approached the interaction---their prompting strategies and ability to diverge from AI suggestions. \citet{fan2024} in the British Journal of Educational Technology coined the term ``metacognitive laziness'' after finding that students using ChatGPT showed significantly fewer metacognitive processes (evaluation, orientation, planning) compared to those working with human experts, even though their short-term output scores were higher.

The most actionable finding comes from \citet{hou2026} in Computers \& Education: in a study with 226 participants, embedding critical thinking training into AI-assisted workflows significantly reduced direct adoption of AI-generated content and produced more creative solutions with higher originality and idea density, even though participants' self-reported critical thinking did not change. Sun, Li, Foo, Zhou, and Lu (2025), in a field experiment randomly assigning employees in a technology consulting firm to receive LLM assistance, found that LLM assistance enhanced creativity especially for employees with high metacognitive strategies---the effect was moderated by the very capacity the B+ Trap threatens. These findings suggest metacognitive habits can be trained faster than metacognitive self-awareness develops---a practical insight for any organization deploying AI creative tools.

The metacognitive escape, however, is inherently limited. It depends on individual human capacity, cannot scale beyond the people who practice it, and works against the grain of tools designed to reward acceptance rather than resistance. This raises a natural question: could an analogous principle---deliberate deviation from the probable---be embedded in the AI architecture itself? The analogy is imperfect: human metacognition involves self-awareness, goal-shifting, and evaluating one's own process, which is qualitatively different from optimizing a different objective function. Nevertheless, the functional similarity is worth exploring: both involve resisting convergence on the most likely output. Before presenting such a framework, it is worth examining what computational creativity theory tells us about the limits of current approaches.

\section{What Current AI Cannot Do: Lessons from Computational Creativity Theory}
Margaret Boden's taxonomy---the dominant theoretical framework in computational creativity---distinguishes three forms of creativity. Combinational creativity produces novel combinations of familiar ideas; exploratory creativity searches within a defined conceptual space; transformational creativity alters the conceptual space itself. Franceschelli and Musolesi (AI \& Society, 2024) mapped these categories to LLMs and found that current systems show capacity for combinational creativity through their vast training data, limited exploratory creativity with careful prompting, but that transformational creativity is ``not achievable through current LLM training solutions'' because RLHF and DPO fine-tuning drive convergence on conventional outputs and LLMs cannot alter their own conceptual spaces.

This theoretical diagnosis aligns precisely with the empirical evidence from Sections 2--4. The B+ Trap is, in Boden's terms, a system that performs adequate combinational creativity (reassembling existing patterns competently) while being structurally incapable of transformational creativity (breaking the patterns themselves). The people who escape the trap are those who supply the transformational component metacognitively---using AI outputs as raw material that they then reshape according to a conceptual framework the AI cannot access. Chen et al.'s (2025) DeepMath-Creative benchmark tested o3-mini, Claude 3.7 Sonnet, Gemini 2.0 Flash, DeepSeek R1, and Qwen QwQ-32B on constructive mathematical problems and found that all models failed to provide effective strategies for open research problems---performance was driven by recombination of memorized patterns rather than authentic creative insight, precisely the combinational-without-transformational limitation Boden's framework predicts.

The SPECS framework \citep{jordanous2012} provides a three-step evaluation procedure grounded in 14 empirically derived creativity components spanning active involvement, dealing with uncertainty, domain competence, originality, and value. Colton's Creative Tripod (2008) requires systems to demonstrate skill (technical proficiency), appreciation (self-evaluation capacity), and imagination (novel output generation). Current LLMs pass the skill criterion easily, show rudimentary appreciation through preference-trained self-evaluation, but fail the imagination criterion in any deep sense---their ``novel'' outputs are recombinations within the statistical boundary of their training distribution, not explorations beyond it.

The Universe of Thoughts framework (2025) directly operationalized Boden's three creativity types for LLMs and found a novelty-utility tradeoff: the most transformative approaches may be difficult to ground in practical solutions, leading models to converge on less transformative but more effective combinational strategies. This is the B+ Trap expressed in theoretical language---models settle for competent recombination because the training objective penalizes the uncertainty inherent in genuine exploration.

\section{The Rebel AI Framework: Training for Productive Deviation}
The exceptional humans described in Section 5 resist convergence cognitively---they critically evaluate, actively diverge, and use AI outputs as points of departure. Each of the following research directions attempts something functionally analogous at the computational level: training systems to deviate from, rather than converge on, their most probable outputs. The parallel is functional rather than mechanistic---optimizing a diversity objective is simpler than human metacognition---but the shared logic of productive deviation connects them. What computational creativity theory in Section 6 identifies as the missing capability---transformational creativity---these approaches begin to address architecturally. No single system yet implements this paradigm in its entirety, but each component is supported by published empirical results.

\subsection{Creative Adversarial Networks: Style Ambiguity as Training Signal}
The most direct approach to adversarial creativity is \citet{elgammal2017}'s Creative Adversarial Network (CAN), which modifies the standard GAN architecture with a deceptively powerful innovation. Where a standard GAN's generator learns only to fool the discriminator into classifying outputs as ``real,'' CAN's discriminator sends two contradictory signals: a standard real/fake classification, and a style ambiguity signal measuring how poorly the discriminator can classify the generated art into known styles (Renaissance, Baroque, Impressionism, etc.). The generator maximizes cross-entropy of the style class posterior with a uniform distribution, receiving steep penalties whenever the discriminator successfully identifies a recognizable style. The theoretical grounding draws on D.E. Berlyne's arousal theory and Colin Martindale's principle that artists increase arousal potential by deviating from established norms.

Trained on 80,000 Western paintings from the 15th--20th century, CAN produced results that human evaluators could not distinguish from contemporary art shown at top art fairs---and actually rated CAN-generated images higher than real contemporary art on several aesthetic dimensions. Without the style ambiguity signal, the same architecture merely reproduces recognizable genres; with it, the generator explores genuinely novel visual territory.

Subsequent work extended adversarial creativity beyond art. \citet{chen2021padgan}'s PaDGAN (ASME Journal of Mechanical Design) introduced a Determinantal Point Process-based loss that simultaneously models diversity and quality, generating airfoil designs with 28\% higher mean quality than vanilla GANs while expanding the design space boundary outside the training distribution. DivGAN (\emph{Artificial Intelligence}, 2023) added a pluggable Diversifier Network inspired by Siamese Networks. For diffusion models, Hierarchical Reward Fine-tuning (HRF, 2024) addresses the observation that reward-based RL fine-tuning causes mode collapse by performing reward-based learning from different diffusion timesteps with an explicit diversity-promoting term.

\subsection{Novelty-Maximizing Rewards: From RLHF to RLHS}
The hypothesis that human surprise or novelty could replace human preference as a reward signal has not been formalized as a distinct paradigm called ``RLHS'' in the published literature. However, several 2025--2026 papers implement closely related ideas. Uniqueness-Aware Reinforcement Learning \citep{hu2026} uses an LLM-based judge to cluster rollouts by high-level solution strategy, then reweights policy advantages inversely with cluster size---correct but rare strategies receive higher rewards. On AIME problems, this achieves up to 100\% recall of canonical solution ideas and sustains higher policy entropy throughout training, preventing exploration collapse without sacrificing pass@1 accuracy.

EVOL-RL \citep{zhou2025evolrl} scores each solution by how different its reasoning is from other concurrently generated responses, combined with majority-voted correctness. Training on label-free data, it lifts Qwen3-4B-Base AIME25 pass@1 from 4.6\% to 16.4\% and pass@16 from 18.5\% to 37.9\%, with generalization beyond math to MMLU-Pro and BBEH. A separate RL framework for creative thinking (2025) derives rewards from Guilford's four dimensions of divergent thinking---novelty, fluency, flexibility, elaboration---finding that creativity-trained models show improved abstraction and flexibility even in code generation and data visualization, suggesting cross-domain transfer of creative capacity.

The most striking result comes from DARLING \citep{li2025darling}, which introduces a learned semantic equivalence classifier to measure diversity beyond surface lexical variation, combined multiplicatively with quality rewards during online RL. On benchmarks spanning instruction following and creative writing at both 8B and 70B scales, DARLING achieves improved Pareto fronts on quality AND diversity simultaneously---directly challenging the assumption that quality and diversity are fundamentally antagonistic. This is the key insight: explicitly optimizing for diversity catalyzes exploration in online RL, which manifests as higher-quality responses. Diverse Preference Optimization (DivPO, Meta AI, \citealt{lanchantin2025}) selects preference pairs based on rarity, yielding 74.6\% greater story diversity while maintaining comparable quality. Creative Preference Optimization (CrPO, \citealt{ismayilzada2025}) injects novelty, diversity, and surprise signals directly into preference optimization using a 200,000+ response dataset, producing a Llama-3.1-8B model that outperforms GPT-4o on both automated and human creativity evaluations---demonstrating that targeted creative fine-tuning of small models can surpass frontier model creativity.

The most theoretically rigorous contribution is Distributional Creative Reasoning \citep{ruizluyten2026}, which proves a ``diversity decay theorem'' showing that STaR, GRPO, and DPO all cause distinct modes of diversity collapse. The framework provides actionable designs ensuring convergence to stable, diverse policies---the first principled theoretical solution to creative mode collapse in language models.

\subsection{The Mathematics of Diversity: Determinantal Point Processes}
Determinantal Point Processes (DPPs)---probability distributions over subsets that assign higher probability to diverse sets---have emerged as the dominant mathematical framework for diversity-promoting training across modalities. The trajectory spans from GDPP (Elfeki et al., ICML 2019), which matched eigenvalues and eigenvectors of real and fake data DPP kernels to alleviate mode collapse in GANs, through PaDGAN's quality-diversity kernel decomposition, to DQO (Diversity Quality Optimization, \citealt{shypula2025}), which constructs a kernel-based similarity matrix over groups of LLM responses per prompt, using the matrix determinant---which measures the squared volume spanned by response embeddings---as a diversity signal. DPP-GRPO \citep{kazimi2025} extends this to text-to-video diffusion models, formulating diversity as a set-level policy optimization problem where the DPP term imposes diminishing returns on redundant samples.

\subsection{Contrastive and Anti-Pattern Methods}
Contrastive Decoding \citep{li2023contrastive} computes the difference between log-likelihoods under a large ``expert'' LM and a small ``amateur'' LM, subject to a plausibility constraint. The insight is that failures like repetition and incoherence are even more prevalent in smaller models, so subtracting the amateur's preferences highlights what the expert uniquely contributes. With zero additional training, contrastive decoding outperforms nucleus sampling, typical sampling, and beam search in MAUVE score across Wikipedia, news, and story domains. SimCTG (Su and Collier, NeurIPS 2022) addresses the root cause---anisotropic token representation distributions---through contrastive training that pushes apart token representations, achieving human-level text generation across 16 languages.

Contrastive Generative Exploration (CGE, 2024) applies contrastive scoring to novelty discovery itself: by computing the contrastive score between a fine-tuned model and its pre-trained base, it generates examples representing novel properties of the fine-tuning data. An iterative version fine-tunes the base on each generated sequence to prevent repetition, forcing systematic exploration of distinct novelties. Verbalized Sampling \citep{zhang2025creative} identifies typicality bias in preference data as a fundamental cause of mode collapse---human annotators systematically favor familiar text, and RLHF amplifies this---then proposes a training-free solution: asking models to generate multiple responses with verbalized probabilities increases diversity by 1.6--2.1$\times$ over direct prompting with no sacrifice in factual accuracy. Min-p sampling (Nguyen et al., ICLR 2025) establishes a minimum probability threshold scaled by the top token's probability, dynamically including diverse options when the model is uncertain; on AlpacaEval Creative Writing, min-p outputs are preferred over top-p, and on GSM8K it achieves higher accuracy AND diversity simultaneously.

A common assumption that temperature scaling alone controls creativity is empirically weak. \citet{peeperkorn2024} found only a weak positive correlation between novelty and temperature and a moderately negative correlation between coherence and temperature in Llama 2-Chat 70B narrative generation evaluated by 36 human participants---far from the ``creativity parameter'' claim. The methods above represent substantially more sophisticated approaches to the diversity problem.

\subsection{Curiosity-Driven Exploration: The Unexplored Frontier}
Pathak et al.'s Intrinsic Curiosity Module (ICM, ICML 2017) established that prediction error in a learned feature space serves as an effective intrinsic reward for exploration---agents trained purely on curiosity crossed significant portions of Super Mario Bros levels without any extrinsic reward. Random Network Distillation \citep{burda2019} simplified this by measuring prediction error against a fixed random network, with novel states producing higher error and thus higher reward.

The critical gap is that neither ICM nor RND has been directly applied to creative text, image, or music generation in published work as of February 2026. The bridge is beginning to form: Motif (2023) elicits preferences from an LLM over observation captions to construct intrinsic rewards, and agents trained only with Motif's intrinsic reward in NetHack surprisingly outperform agents trained on the game score itself. \citet{inch2025} apply novelty search to text-to-image generation via their Wander framework, using an LLM for semantic evolution of prompts and CLIP embeddings to quantify novelty, significantly outperforming existing evolutionary baselines for generating diverse image sets.

A neuroscience finding from \citet{xu2021} adds crucial nuance: novelty and surprise are distinct cognitive mechanisms with separate neural signatures and temporal dynamics. Their SurNoR model shows novelty drives exploration before first reward encounter (ERP responses at 50--100ms), while surprise modulates learning rate (at \textasciitilde{}300ms). This distinction has not yet been operationalized in creative AI systems but suggests that separate novelty and surprise reward channels could yield richer creative exploration than either alone. Combined with Schmidhuber's compression progress theory---creativity as the desire to discover data allowing compression progress, novel but learnable patterns---this points toward a research program where AI creativity emerges from systematic cultivation of productive deviation.

\subsection{Open-Ended Search: The Case Against Objectives}
\citeauthor{stanley2011}'s novelty search \citep{stanley2011} provides the most radical departure from convergent optimization. In a maze navigation task, robots evolved to find novel behaviors---without any fitness-for-goal signal---outperformed robots evolved to reach the goal itself. The insight is counterintuitive but robust: objectives can be deceptive, channeling search toward local optima, while novelty-only search achieves goals as a byproduct of exploring behavioral space. MAP-Elites (\citealt{mouret2015}) extends this by partitioning the behavioral space into a discretized grid and maintaining the highest-performing individual in each cell, simultaneously optimizing for diversity and quality across the entire behavior space---illuminating the search space rather than narrowing it.

Innovation Engines \citep{nguyen2015} combine novelty search with deep learning, evolving images that are simultaneously novel to a trained discriminator and recognized by a deep neural network. The resulting system produced images of remarkable visual complexity and variety---not by being told what to create, but by being rewarded for exploring what had not yet been created. \citet{hughes2024}, in an ICML oral presentation, argued that open-endedness is essential for artificial superhuman intelligence: any fixed objective eventually constrains the system below its potential, while open-ended learning can theoretically continue improving without limit. \citeauthor{stanley2011}'s broader argument, developed in Why Greatness Cannot Be Planned (2015), is that the most significant innovations in both evolution and human culture emerged without predefined objectives---they were stepping stones discovered through exploration rather than targets reached through optimization.

The gap between these open-ended methods and current LLM training remains significant. Standard RLHF is the antithesis of novelty search---it rewards convergence on human-preferred outputs and penalizes deviation from them. The Rebel AI framework is speculative precisely because no system yet integrates adversarial creativity, novelty-maximizing rewards, curiosity-driven exploration, and open-ended search into a unified training paradigm for language models. Each component is proven in its own domain; the integration is the open problem. But the direction is clear: the path out of the B+ Trap requires optimizing for deviation from the probable, not convergence on it.

\section{AlphaEvolve: A Proof of Concept}
Google DeepMind's AlphaEvolve \citep{novikov2025} provides the most compelling existing proof that Rebel AI principles can scale. While not designed as a creative AI system per se, AlphaEvolve implements several Rebel AI components: evolutionary exploration of a solution space (open-ended search), automated evaluation against objective criteria rather than human preferences, and systematic diversification of approaches through population-based methods. The system uses Gemini models as the ``mutation engine'' in an evolutionary algorithm, generating candidate algorithms that are then evaluated against ground truth and selected not just for performance but for diversity of approach.

AlphaEvolve's results are striking. It discovered new algorithms that outperform human-designed solutions for fundamental problems in mathematics and computer science---including finding a new construction for the cap set problem that exceeds the best known human solution. This is transformational creativity in Boden's terms: not recombining known approaches but finding genuinely novel ones that humans had not considered. The system achieves this precisely because it does not optimize for human preferences but for objective performance combined with exploratory diversity.

The limitations of AlphaEvolve as a Rebel AI proof-of-concept are instructive. It operates in domains with clear, automatable evaluation criteria (mathematical proofs, algorithm performance), which does not directly transfer to creative domains where evaluation is inherently subjective. It uses LLMs as components within a broader evolutionary framework rather than training LLMs differently. And it has not been applied to text, image, or music generation. Nevertheless, AlphaEvolve demonstrates that the combination of LLM capability with novelty-seeking search can produce genuinely novel outputs---outputs that would qualify as transformationally creative in domains where such assessment is possible.

\section{Implications for SME Strategy}
While architectures like AlphaEvolve currently remain the domain of specialized research labs, the underlying dynamics of the B+ Trap hold immediate, practical implications for organizations relying on off-the-shelf generative tools. For small and medium enterprises, the challenge is particularly acute. Unlike large corporations that compete on scale, distribution, and brand inertia, SMEs compete disproportionately on creative differentiation: the distinctive voice, the unexpected campaign, the product innovation that carves out a niche. When AI homogenizes creative output across an industry, it erodes precisely the advantage SMEs depend on. This is not speculative---the \citet{liu2025italy} natural experiment shows that AI-homogenized restaurant content was measurably less engaging to consumers, and the Alvero et al. finding that surface-level lexical diversity masks deeper conceptual convergence means SMEs may not even notice when their content has become indistinguishable from competitors'.

First, use AI as amplifier, not author. Reserve AI for drafts, variations, A/B testing, and production tasks while keeping human creative direction over brand voice, emotional storytelling, and strategic positioning. The Doshi and Hauser finding that less creative writers benefit most from AI, while simultaneously anchoring most strongly to AI suggestions, implies that junior team members need explicit protocols for diverging from AI outputs. The Milan restaurant study provides the business case: homogenized content was not just similar---it drove 3.5\% less engagement. For an SME whose marketing budget depends on organic reach, that gap compounds rapidly.

Second, invest in brand infrastructure before AI deployment. Documented voice guidelines, visual standards, and audience personas become both training data for custom configurations and guardrails for human editors. Firms that develop AI configurations internally or customize tools to their specific context achieve more significant returns than those using off-the-shelf defaults: the commodified path is also the homogenized path. SMEs that feed AI tools with distinctive brand assets get outputs that reinforce their differentiation; those that rely on defaults get outputs indistinguishable from every competitor using the same tool.

Third, build explicit co-creation protocols. Cambridge research \citep{luan2025} in Information Systems Research found that human-AI collaboration improves creativity only when participants receive specific instructions on idea co-development; simply providing AI access without structured workflows produces no benefit. The \citet{maier2025} finding that human-led question mode preserves diversity while model-led mode reduces it provides a concrete design principle: structure AI workflows so humans lead the creative direction and AI responds to human queries, not the reverse. This finding undermines the common SME approach of purchasing a subscription and hoping for results. The protocol matters more than the tool.

Fourth, train metacognitive evaluation skills. The \citet{hou2026} intervention study showing that embedded critical thinking training significantly reduced direct adoption of AI content and produced more creative solutions provides a template. The investment is in people, not technology. Given that Fan et al. found ``metacognitive laziness'' sets in rapidly with AI use, and that the CHI timing study found that early AI engagement produces idea fixation while delayed engagement preserves autonomy, ongoing training is essential---not a one-time workshop but a recurring practice embedded in creative workflows.

Fifth, monitor for convergence actively. Regular audits comparing AI-generated content against competitor output can catch similarity drift before it erodes brand differentiation. A growing number of brands are beginning to differentiate based on the deliberate absence of AI in their customer-facing work.

Sixth, demand AI tools designed for divergence, not just quality. The first five recommendations address the human side of the B+ Trap---how teams use AI tools. But the architectural argument of this paper (Sections 6--8) establishes that the tools themselves can be designed differently. Creative Adversarial Networks, novelty-maximizing reward functions such as CrPO and DARLING, Determinantal Point Process-based diversity optimization, and contrastive decoding methods all demonstrate that quality and diversity are not inherently antagonistic---systems can be trained to pursue both simultaneously. As these approaches move from research to product, SMEs should actively evaluate AI tools on diversity metrics, not only output quality. When selecting generative AI vendors, ask whether the system optimizes for distributional diversity or merely for preference alignment. Favor tools that offer diversity-aware generation modes, contrastive or novelty-weighted sampling, or fine-tuning options that preserve output variety. The CrPO result---a fine-tuned 8B model outperforming GPT-4o on creativity evaluations---demonstrates that smaller, diversity-optimized models can surpass larger converge-by-default systems, a finding with direct procurement implications for resource-constrained SMEs. The competitive logic is straightforward: if every company in a sector uses the same convergence-optimized AI tools, creative differentiation becomes a function of who chose different tools, not who used the same tools more skillfully.

For most SMEs, the optimal position lies between full AI adoption and AI avoidance: using AI for operational efficiency while investing disproportionately in the human creative judgment, the diversity-aware tools, and the distinctive brand vision that convergence-optimized language models cannot replicate on their own.

\section{Limitations}
Several limitations should be noted. First, the Rebel AI framework is a synthesis of existing research directions, not an implemented system. Each component has empirical support individually, but no integrated architecture has been tested end-to-end for creative generation. The framework should be read as a research agenda with identified components, not as an engineering blueprint.

Second, the parallel between human metacognition and computational diversity optimization is functional rather than mechanistic. Human metacognition involves self-awareness, goal revision, and reflective judgment---qualitatively richer processes than optimizing a diversity objective function. The claim is that both involve resisting convergence on the most probable output, not that they operate through equivalent mechanisms.

Third, the empirical studies cited vary in methodological rigor, sample size, and ecological validity. The meta-analytic estimates (g = 0.27 for individual creativity, g = $-$0.86 for diversity) aggregate across heterogeneous tasks, populations, and AI systems, and the high heterogeneity (I² = 98.90\% for some comparisons) warrants caution in generalizing effect sizes. Several cited studies are preprints not yet peer-reviewed.

Fourth, the studies documenting the ceiling effect---that top humans outperform AI---do not directly test whether AI pulls strong performers down or whether fewer people attempt to exceed the AI's baseline. The mechanism we propose (motivational thinning rather than capability reduction) is consistent with the evidence but not directly demonstrated by it.

Fifth, the SME implications in Section 9 are derived from the research but are themselves prescriptive recommendations, not empirically validated interventions. The Milan restaurant natural experiment, while suggestive, involves a specific cultural context and content type that may not generalize to all SME sectors.

Sixth, the empirical evidence on creativity spans multiple model generations but has not yet reached the most current frontier. The Creativity Benchmark (Bhat et al., 2025) tested models through O3 and Gemini 2.5 Pro, \citet{haase2025} included o3-mini, and \citet{wright2025} included GPT-5---providing partial evidence that the patterns documented here extend to newer architectures. However, no published study has tested GPT-5.2, Claude Opus 4.6, or Gemini 3 Pro on standardized creativity instruments such as the AUT, TTCT, or creative writing evaluation. Whether these most recent frontier systems exhibit the same creative profile remains an open empirical question. The converging evidence from multiple studies, including the Haase et al. finding of no creative improvement over 18--24 months despite reasoning gains, and the mathematical demonstrations that mode collapse is inherent to standard training objectives (GX-Chen et al., 2025; J. Chen, 2025), suggests the structural logic persists, but dedicated empirical confirmation with the latest models is needed.

\section{Conclusion}
The 2024--2026 research establishes three findings that will shape how organizations approach AI creativity for years to come. First, the B+ Trap is real and structural: RLHF provably narrows output distributions, alignment systematically reduces conceptual diversity, and the resulting creative compression compounds at population scale---the ``Artificial Hivemind'' documented across 70+ models confirms this is an ecosystem-level phenomenon, not a property of any single vendor. Current AI systems excel at combinational creativity but are structurally incapable of transformational creativity---they recombine existing patterns competently but cannot break patterns to create new ones. Second, a key escape mechanism is metacognitive: the people who reach A+ are not using different tools but bringing a fundamentally different cognitive posture---critical evaluation, active divergence, strategic resistance to the probable. Third, the quality-diversity trade-off assumed by standard training is not inevitable: the emerging generation of diversity-aware methods demonstrates that AI systems can be designed to pursue novelty rather than suppress it.

The Rebel AI framework presented here is speculative in its synthesis but empirically grounded in each component. Creative Adversarial Networks, novelty-maximizing rewards, contrastive decoding methods, Determinantal Point Processes, curiosity-driven exploration, and open-ended evolutionary search each represent a tested departure from the convergence logic of standard training. Each mirrors---functionally, if not mechanistically---what exceptional human creators do cognitively: resist gravitational pull toward the median. The parallel is imperfect but productive: both involve deliberately deviating from the most probable output, and the computational versions offer scalability that human metacognition cannot.

For practitioners and business leaders, the message is clear. AI makes everyone a B+ creative. The trap is real, but it is not inevitable---at either the human or the architectural level. The organizations and individuals who thrive will be those who refuse to accept B+ as a ceiling: who invest in the metacognitive skills that let humans escape the trap, who demand AI tools designed for divergence rather than convergence, and who treat creative distinctiveness as a strategic asset worth protecting. The future of creative AI lies not in better approximation of human preferences, but in the systematic cultivation of productive deviation.



\bibliographystyle{plainnat}
\bibliography{references}

\end{document}